% !TeX program = lualatex
% Standalone notebook; compile twice: lualatex 27.tex
% Research batch 39 down to 25, 27 September 2026.
% Original section preserved verbatim except for marked insertion blocks.
\documentclass[11pt,a4paper]{article}
\usepackage[margin=24mm,headheight=14pt]{geometry}
\usepackage{fontspec}
\setmainfont{Latin Modern Roman}
\setsansfont{Latin Modern Sans}
\setmonofont{Latin Modern Mono}
\usepackage{amsmath,amssymb,mathtools}
\usepackage{unicode-math}
\setmathfont{Latin Modern Math}
\usepackage{microtype}
\usepackage{enumitem,xurl,needspace}
\usepackage[dvipsnames]{xcolor}
\usepackage{hyperref}
\hypersetup{unicode=true,colorlinks=true,linkcolor=MidnightBlue,urlcolor=MidnightBlue,
 pdftitle={Research notebook: Problem 27},
 pdfauthor={Open Problems Math research notebook},bookmarksnumbered=true}
\usepackage{bookmark,titlesec,fancyhdr}
\titleformat{\section}{\Large\bfseries\raggedright}{\thesection}{0.7em}{}
\pagestyle{fancy}\fancyhf{}
\fancyhead[L]{\small\sffamily Open Problems Math\enspace|\enspace Research notebook}
\fancyhead[R]{\small\sffamily Problem 27}
\fancyfoot[L]{\small\sffamily 27 September 2026}
\fancyfoot[R]{\small\sffamily \thepage}
\setlength{\parindent}{0pt}
\setlength{\parskip}{5pt plus 1pt}
\setlength{\emergencystretch}{3em}
\setcounter{secnumdepth}{1}
\setcounter{tocdepth}{1}
\setlist[itemize]{leftmargin=3em,itemsep=5pt,topsep=4pt,parsep=0pt}
\allowdisplaybreaks
\newcommand{\N}{\mathbb{N}}
\newcommand{\Z}{\mathbb{Z}}
\newcommand{\Q}{\mathbb{Q}}
\newcommand{\R}{\mathbb{R}}
\newcommand{\C}{\mathbb{C}}
\newcommand{\F}{\mathbb{F}}
\newcommand{\A}{\mathbb{A}}
\newcommand{\PP}{\mathbb P}
\newcommand{\E}{\mathbb{E}}
\newcommand{\eps}{\varepsilon}
\newcommand{\dd}{\,\mathrm d}
\newcommand{\sref}[2]{\hyperlink{source:#1:#2}{[S#2]}}
\newcommand{\eref}[2]{\hyperlink{extra:#1:#2}{[E#2]}}
\newif\ifshowcatalogue
\showcataloguetrue
\newcommand{\cataloguescope}[1]{\ifshowcatalogue
 \paragraph{Exact scope in the supplied catalogue.}
 \begin{quote}\small #1\end{quote}\fi}
\numberwithin{equation}{section}
\begin{document}
\setcounter{section}{26}
% BEGIN PRESERVED SECTION
% ================================================================
% problemId: problem.p-versus-bpp-does-bpp-p
% source releaseRank: 27; source releaseStatus: open
% exactTarget SHA-256: 4dc5eb5c4530c8af4159767a2e6e06ac59a75b0bc2bb55658d4bf306465dfb54
\clearpage
\section{\texorpdfstring{P versus BPP (Does BPP = P?)}{P versus BPP (Does BPP = P?)}}\label{27}
\subsection{Definitions and mathematical statement}
A language $L$ is in $\mathsf{BPP}$ if a randomized polynomial-time algorithm $A$ satisfies $\Pr[A(x)=1_L(x)]\ge2/3$ for every input $x$. Membership in $\mathsf P$ requires a deterministic polynomial-time decision algorithm with no errors. Determine whether
\[
 \forall L\in\mathsf{BPP},\quad L\in\mathsf P.
\]
The simulation must work in the ordinary uniform model, not merely through nonuniform advice or for most inputs.
\par\smallskip\textit{Formulation and concept references:} \sref{27}{1}.
\cataloguescope{Determine whether every language computable by a probabilistic polynomial-time algorithm with bounded two-sided error (BPP) can be computed by a deterministic polynomial-time algorithm, i.e., whether BPP = P.}
\subsection{Short English statement}
Can randomness always be removed from an efficient bounded-error decision algorithm without losing polynomial-time efficiency?
% BEGIN RESEARCH 27
\subsection{Research attempt}

\paragraph{Current frontier.}
A central established conditional result is that exponential circuit
hardness in $\mathsf E$ at all sufficiently large lengths suffices for
$\mathsf{BPP}=\mathsf P$. The all-length and nonuniform-hardness
qualifications cannot be removed from that implication. Vadhan gives
this in Corollary~7.64 and discusses the infinitely-often distinction
immediately afterward. \sref{27}{1}. Doron--Moshkovitz--Oh--Zuckerman's
May 2026 manuscript improves generator and hitting-set parameters for
algorithms satisfying specified insensitivity properties, within the
hardness-versus-randomness program; it does not establish unconditional
derandomization of arbitrary BPP machines. \sref{27}{3}, \eref{27}{1}.
The paper's abstract and parameter scope, and the recent course account
\sref{27}{2}, were checked. The argument below does not assume the hard
function or the insensitivity property for every machine.

\paragraph{Attack route.}
The probabilistic method gives a short list of random tapes that is good
for every input of a fixed length. The missing step is generating such a
list uniformly. I combine this finite existence proof with a robust
conditional-expectation selection rule, asking exactly what an efficiently
computable pessimistic quantity would have to provide. The competing
approach replaces the tapes with a bounded-independence distribution;
an explicit BPP acceptance predicate tests whether that alone suffices.
Nonuniform derandomization and conditional expectations are classical
techniques treated separately in \sref{27}{1}, Sections~3.2 and~3.4;
the following calculation puts their computational gap in one place.

\paragraph{Attempt.}
Fix a BPP machine $A$ for a total language $L$. On inputs of length $n$,
pad its randomness to exactly $r(n)\ge1$ bits, with $r(n)$ polynomially
bounded and computable. Every input $x$ has error probability at most
$1/3$. Take $k=24(n+2)+1$ independent tapes $R_1,\ldots,R_k$, so $k$ is
odd, and let $A^{[k]}$ return their majority vote.

\textbf{Lemma 27.1 (a simultaneously good polynomial-length tape list).}
There is a fixed tuple $s_n=(R_1,\ldots,R_k)$ such that
$A^{[k]}(x;s_n)=1_L(x)$ for every $x\in\{0,1\}^n$.

\emph{Proof.} For fixed $x$, let $S$ count erroneous independent runs.
Since each has error probability at most $1/3$,
$\mathbb E2^S\le(4/3)^k$. Markov's inequality gives
\[
 \Pr[S\ge k/2]\le\left(\frac{2\sqrt2}{3}\right)^k=\rho^k.
\]
As $k$ is odd, majority error is included in this event. Also
$\rho^2=8/9$ and $(8/9)^{12}<1/4$ (an exact rational inequality).
The expected number of inputs on which the majority errs is at most
\begin{equation}\label{r27:smallpotential}
       2^n\rho^k<2^n4^{-(n+2)}\le1/16.
\end{equation}
If every tuple were bad on some input, that integer-valued count would
always be at least one. Hence a tuple with no errors exists.\hfill$\square$

The tuple has $kr(n)=\operatorname{poly}(n)$ bits. Hardwiring one such
tuple for each $n$ proves a nonuniform circuit simulation, not the desired
uniform algorithm. The proof has made no choice rule for $s_n$.

\textbf{Lemma 27.2 (a robust conditional-expectation criterion).}
Write $B=kr(n)$ and regard the tape tuple as a uniform $B$-bit string
$S$. Let
\[
 Z(S)=\#\{x\in\{0,1\}^n:A^{[k]}(x;S)\ne1_L(x)\},
 \qquad W(p)=\mathbb E[Z(S)\mid S\text{ extends }p].
\]
Suppose that, for this fixed machine, a uniform deterministic
polynomial-time procedure computes on every prefix $p$ a rational number
$\widetilde W(p)$ such that
\begin{equation}\label{r27:estimate}
             |\widetilde W(p)-W(p)|\le\frac1{8B}.
\end{equation}
Then $L\in\mathsf P$.

\emph{Proof.} Starting at the empty prefix, compute the estimates for
$p0,p1$ and append the bit with smaller estimated value. Since
$W(p)=(W(p0)+W(p1))/2$, at least one true child value is at most $W(p)$.
The chosen true value is at most this smaller true value plus $2/(8B)$.
After $B$ choices the resulting string $s_n$ therefore satisfies
\[
          Z(s_n)=W(s_n)\le W(\varnothing)+1/4<1.
\]
The nonnegative integer $Z(s_n)$ must be zero. Generating the tuple takes
$2B$ polynomial-time estimation calls; the majority of the $k$ runs is
also polynomial time. The construction is uniform at every length,
including small lengths, and uses no advice.\hfill$\square$

This is a genuine sufficient algorithmic condition, but not an
independently established estimator. In particular, $W$ is defined using
the unknown correct language bits. On a full string it counts errors over
$2^n$ inputs. Rewriting it as a finite sum does not make it polynomial-time
computable. Replacing the language bits by majority over \emph{all}
random tapes merely inserts further exponentially large sums. The
conditional-expectation identity proves an advantageous child exists;
it does not by itself identify that child efficiently.

The absolute accuracy in \eqref{r27:estimate} is only inverse-polynomial
for the unnormalized count $W$, but $W$ sums over $2^n$ inputs. An
inverse-polynomial error for the \emph{normalized} average $2^{-n}W$
is far too coarse for the same selection proof: after rescaling it could
hide exponentially many bad inputs. This is another precise difference
between an almost-everywhere simulation and the total-language target.

\medskip
\textbf{Testing whether bounded independence repairs the selection gap.}
Let $r\ge2$ and let $D$ be uniform on the even-parity strings
$u\in\{0,1\}^r$, with an additional independent uniform bit $b$.
Any $r-1$ coordinates among the $u$ bits are uniform and independent:
once they are fixed, exactly one last bit satisfies the parity equation.
Consequently $(u,b)$ is $(r-1)$-wise independent. Nevertheless, the
polynomial-time acceptance predicate
\begin{equation}\label{r27:parity}
                 A_*(u,b)=\operatorname{Parity}(u)\ \lor\ b
\end{equation}
has probability $3/4$ under fully uniform bits and probability $1/2$
under $D$. As a machine ignoring its ordinary input, $A_*$ is a valid
BPP decider for the all-yes language with error $1/4$. The bounded-
independence replacement loses the required acceptance margin, even
though its independence order grows linearly with the tape length.
The companion script verifies the marginals and the two probabilities
exactly for small $r$.

This does not prove a separation: the all-yes language has a trivial
deterministic decider. Nor does it refute bounded-independence generators
for restricted circuit classes. It refutes the specific proposed
universal inference that high independence of random coordinates alone
preserves every BPP acceptance gap. Enumeration of that particular
support and a tie rule is not a general replacement for a proof of
pseudorandomness against the algorithm being simulated.

\medskip
\textbf{A second sufficient condition, with its seed length exposed.}
Suppose instead that for every BPP machine $A$ there is a uniform
polynomial-time computable map
\[
 G_{A,n}:\{0,1\}^{d(n)}\longrightarrow\{0,1\}^{r(n)},
             \qquad d(n)\le C_A\log_2(n+2),
\]
whose output changes $A$'s acceptance probability by less than $1/12$
on every input of length $n$. Enumerating all seeds then computes an
acceptance fraction that is at least $7/12$ on yes-inputs and at most
$5/12$ on no-inputs. Comparison with $1/2$ decides $L$ deterministically
in polynomial time. This elementary conversion is the reason the
logarithmic seed length and uniform generator evaluation matter.
\sref{27}{1}, Chapters~3 and~7. A polynomial-length seed yields an
exponential enumeration, and a different hardwired generator for each
length yields only advice unless uniformly generated.

\paragraph{Stress test.}
The existence proof and both conditional simulations handle a total
language, not a promise on a favorable set of inputs. They permit the
polynomial exponent to depend on $A$ but not on the individual input.
The amplification probability uses independent runs of $A$; it is not
applied to the parity-conditioned tapes as though they were independent.
Lemma~27.2 explicitly budgets every approximation error, so no exact
real-arithmetic oracle is hidden in the final selection rule. Its
hypothetical estimator must output ordinary polynomial-bit rationals.

The hard-function theorem in the frontier would imply the generator
condition using much deeper machinery, but its circuit lower bound is
not proved here. The insensitivity result also has additional hypotheses,
not a theorem that every BPP machine meets them. \eref{27}{1}.
Neither the small good list nor failure of the parity replacement gives
an unconditional class separation. The companion finite tests check the
amplification constants and the parity example only; a numerical test
cannot verify all BPP machines.

\Needspace{8\baselineskip}
\paragraph{Outcome.}
\textbf{Outcome: Conditional reduction.}

The attempt proves a uniform derandomization criterion based on an
accurate conditional-error-count estimator, including a quantitative
error budget, and records the short-generator alternative with every
parameter explicit. It also demonstrates why nonconstructive short lists
and high-order independence do not themselves meet these conditions.
No polynomial-time estimator or general short generator is constructed
unconditionally. The remaining gap is computational, not the existence
of a good tuple. The constituent techniques are classical, and no
historical novelty or complete resolution of $\mathsf P=\mathsf{BPP}$
is claimed.

% END RESEARCH 27

\subsection{Sources}
\begin{itemize}\raggedright
\item[\textnormal{[S1]}]\hypertarget{source:27:1}{} S. Vadhan, 'Pseudorandomness', Foundations and Trends in Theoretical Computer Science 7(1-3):1-336, 2012.
\url{https://people.seas.harvard.edu/~salil/pseudorandomness/pseudorandomness-Aug12.pdf}.
{\small\textit{Catalogue formulation source.}}
\item[\textnormal{[S2]}]\hypertarget{source:27:2}{} CS 6810: Theory of Computing, Lecture 15, March 17, 2026.
\url{https://www.cs.cornell.edu/courses/cs6810/2026sp/lec15.pdf}.
{\small\textit{Catalogue research/status reference; not a proof endorsement.}}
\item[\textnormal{[S3]}]\hypertarget{source:27:3}{} Pseudorandomness Beating the Hybrid Argument for Insensitive Algorithms.
\url{https://eccc.weizmann.ac.il/report/2026/082/}.
{\small\textit{Catalogue research/status reference; not a proof endorsement.}}
% BEGIN ADDED REFERENCES 27
\item[\textnormal{[E1]}]\hypertarget{extra:27:1}{} Dean Doron, Dana Moshkovitz,
Justin Oh and David Zuckerman, \emph{Pseudorandomness Beating the Hybrid
Argument for Insensitive Algorithms}, ECCC TR26-082, revision~1,
17 May 2026.
\url{https://eccc.weizmann.ac.il/report/2026/082/}.
{\small\textit{Version-specific supplement to [S3]. Abstract, generator
parameters, and extra insensitivity hypotheses consulted 27 September 2026;
not an unconditional BPP derandomization theorem.}}

% END ADDED REFERENCES 27
\end{itemize}

% END PRESERVED SECTION
\end{document}
